\subsection{Completion length and warm serial timing}
\label{sec:serial_tokens}
Table~\ref{tab:serial_tokens} reports the same 100 balanced records
per mode and fit (20 per class), including end-of-turn tokens.
All recorded output-limit counts are zero. These are serial benchmark
outputs; their token counts are not substituted for full-test metrics.
Adapter execution and response length both change between modes,
so this table describes their joint observation rather than isolating
either contribution to latency.
\begin{table}[tb]\centering\small
\caption{Per-fit serial completion length and latency. Token counts and
latency are means over 100 requests in each mode; seconds include prompt
rendering, tokenisation and synchronised generation.}
\label{tab:serial_tokens}
\begin{tabular}{llrrrr}\toprule
Model & Fit seed & Base tokens & Adapted tokens & Base (s) & Adapted (s) \\\midrule
Gemma 4 & 1 & 9.48 & 9.57 & 0.475 & 0.647 \\
Gemma 4 & 2 & 9.48 & 9.60 & 0.482 & 0.659 \\
Gemma 4 & 3 & 9.48 & 9.54 & 0.490 & 0.661 \\
Qwen3 & 1 & 9.27 & 9.35 & 0.421 & 0.539 \\
Qwen3 & 2 & 9.27 & 9.30 & 0.422 & 0.539 \\
Qwen3 & 3 & 9.27 & 9.34 & 0.421 & 0.539 \\
Phi-4 & 1 & 13.01 & 9.36 & 0.487 & 0.435 \\
Phi-4 & 2 & 13.01 & 9.35 & 0.487 & 0.435 \\
Phi-4 & 3 & 13.01 & 9.32 & 0.490 & 0.440 \\
\bottomrule\end{tabular}\end{table}
